\chapter{Bilateral original-source reductions, maximal integral lifts,
and subquadratic support bounds}
{\small\sffamily Source: \nolinkurl{repository/collatz_reconstruction/research_program/bilateral_root_bounds_20260916/note.md}.\par}


16 September 2026. Research continuation in
\nolinkurl{KokunoYumeto/collatz-workbench}. The authenticated base is
\nolinkurl{b2dac4c2cc4a5eedde0af36d6749390c77fa4a62},
\nolinkurl{ternary_join_reduction_20260916}. The original integral
receiver is \nolinkurl{intrinsic_zero_firstjet_20260915}. All earlier
source files remain unchanged.

This note proves two complete new arithmetic families at every positive
ternary valuation, their exact disjointness from the preceding layer
families, and an all-depth refinement that selects the smallest source
within each admitted leading-run template. It also gives a complete
finite two-sided catalogue and a subquadratic bound on the original
vertex support of the whole-source retraction. The first-crossing window
is verified through 11,610 odd steps by a complete original-source
certificate and an independent all-row audit.

These are bounds on specified constructions and proved source domains.
They do not assert that the residual root set is empty, that every orbit
crosses, or that a possible positive nontrivial cycle or nonperiodic
component has been found. No estimate from the Zeta programme is
imported. The general arithmetic and algebraic arguments below use no
theorem on linear forms in logarithms.

\section{1. Fixed original coordinates and the comparison receiver}

Write X for the positive odd integers and retain

\SourceMath{T(n)=\frac{3n+1}{2^{a(n)}},\qquad a(n)=\nu_2(3n+1).}{1}

An exponent word w=(a\_1,\ldots,a\_k) has original prefix sums A\_j,
total A, and

\SourceMath{L=3^k,\quad U=2^A,\quad
 C=\sum_{j=0}^{k-1}3^{k-1-j}2^{A_j},\quad
 f_w(n)=(Ln+C)/U.}{2}

The empty word has (L,U,C)=(1,1,0). Its complete positive odd source is
1+2t. For every word retain the original cylinder and image

\SourceMath{\rho+2Ut\longmapsto\eta+2Lt,\qquad t\ge0,
\quad \rho\equiv(U-C)L^{-1}\pmod{2U},\quad0<\rho<2U,
\quad\eta=(L\rho+C)/U.}{3}

Their inverse parameters are (n-rho)/(2U) and (y-eta)/(2L) on the
displayed images. The affine numerator satisfies the chronological
recurrence C'=3C+U. The terminal congruence in (3) gives every prefix
oddness: reduce the full numerator modulo 2\^{}\{A\_j+1\}, retain the
j-prefix term, the next term of exact valuation A\_j, and the later
terms divisible by 2\^{}\{A\_j+1\}. The resulting prefix numerator is
2\^{}\{A\_j\} modulo 2\^{}\{A\_j+1\}. The successive original equations
then force the exact valuations a\_j. This proves the word/cylinder
correspondence without replacing exact valuations by divisibility alone.

In particular 0\textless eta\textless2L. Positivity follows from
rho\textgreater0. The same congruences at the specified negative source
rho-2U give the integral odd terminal eta-2L. Every original odd-return
image of a negative odd integer is negative; hence eta-2L\textless0.
This proves the image-anchor bound by two evaluations of the same affine
map. No identification of the positive and negative source domains is
made.

The absolute first-jet complex has original edge and vertex bases
E\_n,V\_n, including E\_1,V\_1. Put

\SourceMath{\mathbb D=\mathbb Z[\epsilon]/(\epsilon^2),\quad q=1+\epsilon,
 \qquad d_\epsilon E_n=V_n-qV_{T(n)}.}{4}

For every integer k, q\^{}k=1+k epsilon; q\^{}\{-1\}=1-epsilon is an
integral unit. Neither a cycle period nor a nonunit integer is inverted.
The ordinary incidence is d=J-P. The specified comparison kernel is

\SourceMath{\mathfrak B_X=\ker\bigl(H^1(K_\epsilon)\to H^1(K)\bigr),\qquad
 \xi_n=[\epsilon V_n].}{5}

For an original finite path l:n-\textgreater y of length s let B\_l(n)
be its weighted chain, with coefficient q\^{}j on its j-th edge.
Telescoping gives

\SourceMath{d_\epsilon B_l(n)=V_n-q^sV_y.}{6}

Thus two ORIGINAL paths l:n-\textgreater y and r:m-\textgreater y have
the exact comparison

\SourceMath{\boxed{d_\epsilon\bigl(B_l(n)-q^{s-t}B_r(m)\bigr)
       =V_n-q^{s-t}V_m,\quad t=|r|.}}{7}

Multiplication by epsilon gives xi\_n=xi\_m. This equality does not by
itself make either original class zero. It identifies their images
through the original finite integral relation, with both paths and their
source labels retained. In particular the integer m below is a smaller
source with a common future, not necessarily a forward iterate of n.

The predecessor proves that the full kernel (5) is a direct sum of
Z/m\_C Z for actual nontrivial cycle components and Z for actual
nonperiodic components. Its finite integral witness equations are db=0,
dc-Pb=V\_n.~They imply an actual path to 1, extracted from the original
support, since the cycle-component augmentation would otherwise give
-k*m\_C=1 with m\_C\textgreater1 or with no cycle. We retain that exact
receiver and independently replay the supplied original validator; no
quotient coefficient is silently rationalized.

\section{2. Two nontrivial legs on every ternary layer}

Fix b\textgreater=1 and e in \{2,6\}. Put

\SourceMath{h_e=(2^e+2)/3\in\{2,22\},\qquad
 a=a(b,e)=\min\{a\ge1:2^{a+e+2b-1}<3^{a+b}\},}{}
\SourceMath{J=2^{e+2b-1},\quad Q=3^{a+b},\quad K=2^aJ<Q.}{8}

These definitions involve only integer comparisons. The minimum exists,
with

\SourceMath{a(b,e)\le b+2e-2,\qquad
 a(b+1,e)-a(b,e)\in\{0,1\}.}{9}

At a=b+2e-2 the ratio K/Q is (8/9)\^{}\{b+e-1\}\textless1. Increasing b
multiplies the ratio by 4/3; increasing a once more multiplies it by
2/3. Their product is 8/9\textless1, which proves the upper increment,
and monotonicity proves the lower increment.

The left ORIGINAL word is l=(1,2,1), whose map is (27n+23)/16 and whose
source is n=27 modulo32. The right word is

\SourceMath{r_{b,e}=(1)^a,(e),(2)^{b-1},(1),(1),(3).}{10}

It has a+b+3 odd returns and exponent sum a+e+2b+3. Define the exact
source set

\SourceMath{\mathcal S_{b,e}=\{n>0:n\equiv27\pmod{32},\quad
                         Jn+h_e3^b\equiv0\pmod Q\}.}{11}

The moduli 32 and Q are coprime, so this is one complete positive
residue class modulo32Q. Let its representative in (0,32Q) be n\_0 and
put

\SourceMath{m_0=(Kn_0+2^ah_e3^b)/Q-1,\qquad y_0=(27n_0+23)/16.}{}

\subsection{Theorem 1. Complete bilateral joins at every positive
ternary valuation}

For every b,e above, the entire designated original source fibre is

\SourceMath{\boxed{\begin{aligned}
 n(t)&=n_0+32Qt,\\
 m(t)&=m_0+32Kt,\\
 y(t)&=y_0+54Qt,\qquad t\in\mathbb Z_{\ge0},
 \end{aligned}}}{12}

with the original paths

\SourceMath{n(t)\xrightarrow{\ (1,2,1)\ }y(t)
       \xleftarrow{\ r_{b,e}\ }m(t),\qquad0<m(t)<n(t).}{13}

Every n(t) has actual nu\_3(n(t))=b. Its unit n(t)/3\^{}b is 2 modulo3
for e=2 and 1 modulo3 for e=6. All these sources have no incoming
original odd-return edge.

\textbf{Proof with all intermediate values.} Congruence (11) forces
3\^{}b\textbar n.~Write n=3\^{}b z. Then Jz+h\_e=0 modulo3\^{}a. Since
J=2 modulo3 and h\_e is a ternary unit, z is a unit of the asserted
class. Thus b is the exact valuation. Set W=(Jz+h\_e)/3\^{}a. It is a
positive even integer of exact dyadic valuation one: J is divisible by
eight and h\_e=2 modulo4, while the divisor is odd. Define m=2\^{}aW-1,
which is positive odd. Its first a original values are

\SourceMath{x_j=3^j2^{a-j}W-1\quad(0\le j\le a).}{14}

For j\textless a, x\_j+1 is divisible by four, giving exact next
exponent one. At j=a,

\SourceMath{x_a=Jz+h_e-1=2^{e-1}(4/3)^b n+h_e-1.}{15}

Because 3h\_e-2=2\^{}e, the next exponent is exactly e and its image is

\SourceMath{v_0=1+2\,4^{b-1}n/3^{b-1}.}{}

The following b-1 exponents are exactly two, with original values

\SourceMath{v_j=1+2\,4^{b-1-j}n/3^{b-1-j}\quad(0\le j\le b-1).}{16}

For j\textless b-1 the nonconstant part is divisible by eight, which
proves exactness of those valuations. Their last value is 2n+1. The
remaining original values are

\SourceMath{2n+1\xrightarrow{1}3n+2\xrightarrow{1}(9n+7)/2
                \xrightarrow{3}(27n+23)/16.}{17}

The residue n=27 modulo32 proves all three exact valuations. It also
proves the left path n -\textgreater{} (3n+1)/2 -\textgreater{} (9n+5)/8
-\textgreater{} (27n+23)/16 with exponents (1,2,1). Conversely equality
of the original two endpoint formulas gives m=(Kn+2\^{}ah\_e3\^{}b)/Q-1
and hence (11); the left path requires n=27 modulo32. Thus the source
domain is exact, rather than a subset justified only by samples.

The height difference is controlled by

\SourceMath{m=(K/Q)n+h_e(2/3)^a-1.}{18}

The defining comparison K/Q\textless1 gives

\SourceMath{h_e(2/3)^a<(2h_e/2^e)(3/4)^b<1.}{19}

Both terms in (18) therefore make m\textless n.~Positivity follows from
m=2\^{}aW-1. CRT gives (12); subtraction and division in any one of the
three displayed progressions supplies the inverse parameter. An incoming
edge to n would require 3x+1=2\^{}c n=0 modulo3, which is impossible.
This proves every assertion. QED.

The smallest displayed case b=1,e=2 has a=3:

\SourceMath{\boxed{1275+2592t\longrightarrow2153+4374t
                  \longleftarrow1007+2048t.}}{20}

At t=0 the original left path is 1275,1913,1435,2153. The original right
path is \nolinkurl{1007,1511,2267,3401,2551,3827,5741,2153.} The
difference of the two sources is 268+544t\textgreater0. This family can
already overlap a finite catalogue; its full integral relation remains
valid there.

The case b=2,e=6 has a=10:

\SourceMath{n(t)=8808219+17006112t,\qquad
 m(t)=8689663+16777216t,}{21}

with the common endpoint (27n(t)+23)/16. This is an unbounded family of
original sources beyond the finite verification interval below, not an
extrapolation from that interval.

\section{3. The old and new layer families have no common source}

The preceding published families use

\SourceMath{J^{\rm old}_{b,e}=2^{e+2b},\quad
 J^{\rm old}_{b,e}(n/3^b)+h_e=0\pmod{3^{a^{\rm old}(b,e)}},
 \quad n=3\pmod4,}{22}

where a\^{}old is the least positive a with
2\^{}\{a+e+2b\}\textless3\^{}\{a+b\}. Their unit classes are 1 modulo3
for e=2 and 2 modulo3 for e=6.

\subsection{Theorem 2. Exact disjointness, including the first
nonreduced lift}

The union of the new families (11) is disjoint from the union of all the
old families (22). Families with different b are disjoint by the
original valuation; at a fixed b the two e values are disjoint by the
original unit class.

\textbf{Proof.} The only possible mixed intersections at a fixed b pair
new e=2 with old e=6, or new e=6 with old e=2. In the first case
J\^{}old\_6=32J\^{}new\_2. Multiplying the new congruence by 32 and
subtracting the old gives the constant

\SourceMath{32h_2-h_6=64-22=42,\qquad\nu_3(42)=1.}{23}

In the second case J\^{}new\_6=8J\^{}old\_2, giving

\SourceMath{h_6-8h_2=22-16=6,\qquad\nu_3(6)=1.}{24}

Every precision in these possible pairs is at least two; this follows
either from (8),(22) at b=1 and monotonicity, or by checking that a=1
cannot contract. Either intersection would therefore force 9\textbar42
or 9\textbar6, a contradiction. QED.

The matching congruences can agree modulo3 and disagree modulo9.
Equations (23)-(24) retain exactly the obstruction at the next lift.
Thus matching a unit class alone would not establish an overlap; the
original coefficient maps and prime-power precision are essential. This
statement concerns the displayed integer source congruences, not an
unproved identification with a Zeta conductor.

\section{4. Maximal leading runs: the whole extra-depth axis in one
finite calculation}

The original source is preserved while all allowed leading runs are
compared. For n in (11), set

\SourceMath{Z=J(n/3^b)+h_e,\qquad A_{\max}=\nu_3(Z).}{25}

\subsection{Theorem 3. Exact maximal lift and its full parameter strata}

The possible original leading runs in this fixed template have
a'\textless=A\_max. Every contracting one has
a\textless=a'\textless=A\_max, with smaller source

\SourceMath{m_{a'}=2^{a'}Z/3^{a'}-1.}{26}

The smallest of these sources is m\_\{A\_max\}. For a'\textless A\_max,
the original edge T(m\_\{a'+1\})=m\_\{a'\} has exact exponent one. The
original maximum vertex on these right paths is independent of a': their
initial rises end at the same Z-1 and all subsequent vertices coincide.

\textbf{Proof.} Integrality of the forced endpoint formula requires
3\^{}\{a'\}\textbar Z, hence exactly a'\textless=A\_max. Division by an
odd power leaves nu\_2(Z/3\^{}\{a'\})=1, so (14)-(17) prove the actual
words. All a'\textgreater=a contract by the comparison in (8). Equation
m\_\{a'+1\}+1=(2/3)(m\_\{a'\}+1) proves strict decrease and, since the
first odd value has its plus-one valuation at least two, the exact
exponent-one edge. The initial rise ends at Z-1 for every a', and is
increasing up to that point; the common remaining path proves the
maximum assertion. QED.

This theorem exhausts the admitted leading-run axis for the specified
template. It makes no claim that every possible right-word shape is one
of these templates.

There is an exact all-depth partition of the ORIGINAL family parameter.
Write n=n\_0+32Qt and

\SourceMath{W_0=(J(n_0/3^b)+h_e)/3^a.}{}

Then

\SourceMath{Z=3^a(W_0+32Jt).}{}

For an extra depth r\textgreater=0 and a retained unit u in \{1,2\}, its
entire stratum is

\SourceMath{\boxed{t=t_0+3^{r+1}s,\quad
 t_0\equiv(3^r u-W_0)(32J)^{-1}\pmod{3^{r+1}},
 \quad0\le t_0<3^{r+1},\quad s\ge0.}}{27}

Here A\_max=a+r. The complete source and smaller-source steps are
respectively

\SourceMath{32Q\,3^{r+1},\qquad 2^{a+r}(32J)3.}{28}

Their constants are obtained from n(t\_0) and
2\^{}\{a+r\}(W\_0+32Jt\_0)/3\^{}r-1. Both are original positive
integers. The inverse is s=(t-t\_0)/3\^{}\{r+1\}; exact valuation r
follows from the displayed nonzero unit u. Every t\textgreater=0 belongs
to exactly one such stratum because W\_0+32Jt is positive. The same
inverse parameters compute its common original endpoint. Thus arbitrary
extra depth is retained without a truncation or a generic large-depth
flag.

\subsection{Exact comparison with the minimal-run relation}

Let B\_min be (7) for the minimal-run pair, with clock k\_min. Let I be
the original r-edge incoming path from m\_\{A\_max\} to m\_a. Its
weighted chain obeys d\_epsilon I=V\_\{m\_Amax\}-q\^{}r V\_\{m\_a\}. The
maximal pair therefore has chain

\SourceMath{\boxed{B_{\max}=B_{\min}-q^{k_{\min}-r}I.}}{29}

The two original right words concatenate as 1\^{}r followed by the
minimal right word, so this is an equality of original weighted chains,
not only of boundaries. Maximization does not supply a different
root-exclusion domain: it refines the same admitted family by composing
its original relation with r existing incoming ones. That distinction is
retained in the implementation.

There is also a source-size bound on the chosen run. Put k=floor(log\_3
n), computed as the largest integer with 3\^{}k\textless=n.~Since
J+h\_e\textless3\^{}\{e+2b-1\} and n/3\^{}b\textless3\^{}\{k+1-b\},

\SourceMath{Z\le(J+h_e)n/3^b<3^{k+b+e},\qquad
 A_{\max}\le k+b+e-1.}{30}

This is a proved finite search bound and an exact one-step valuation
evaluation; it is not an assumption about an arbitrary orbit's stopping
time.

\section{5. Complete two-leg catalogue with both nonempty original
words}

For original words l,r, keep all five affine data (L,U,C,rho,eta) from
(2)-(3). Put g=gcd(L\_l,L\_r). Their target images intersect exactly
when

\SourceMath{\eta_l\equiv\eta_r\pmod{2g}.}{31}

Let y\_0 be the least common image value at least max(eta\_l,eta\_r),
and put u\_0=(y\_0-eta\_l)/(2L\_l), v\_0=(y\_0-eta\_r)/(2L\_r). Then all
compatible pairs are

\SourceMath{\begin{aligned}
 n(t)&=n_0+2U_l(L_r/g)t,&n_0&=\rho_l+2U_lu_0,\\
 m(t)&=m_0+2U_r(L_l/g)t,&m_0&=\rho_r+2U_rv_0,\\
 y(t)&=y_0+2(L_lL_r/g)t,&t&\ge0.
\end{aligned}}{32}

The common target and its period give the inverse t. The original source
parameters are u=u\_0+(L\_r/g)t and v=v\_0+(L\_l/g)t, with the inverse
obtained by subtraction and division. The height condition is exactly

\SourceMath{(n_0-m_0)+\bigl(2U_lL_r/g-2U_rL_l/g\bigr)t>0.}{33}

It specifies an integer interval with either slope sign, which is
retained even when empty. Formula (31) follows from the congruence
compatibility criterion; (32) follows by adding the common target
period. Thus both inverse maps and all images are proved. No
independent-distribution premise is used.

The new finite domain has these five left words:

\SourceMath{(1,1),\ (1,2),\ (1,1,1),\ (1,1,2),\ (1,2,1).}{34}

They are all the indicated length-two and length-three words before
coefficient crossing. For each left word and each right length
1\textless=j\textless=12, include EVERY positive exponent word r with

\SourceMath{2^{A_r}L_l<U_l3^j.}{35}

The total sum A\_r has a finite exact ceiling B(l,j), computed by
integer powers. There are binom(B(l,j),j) right words when
B(l,j)\textgreater=j. This is the positive- composition identity
obtained by adjoining the unused budget as a final nonnegative
coordinate. Hence the complete domain has 112,050 pairs.

Of these, 24,107 satisfy (31). The complete certificate proves that each
of these has domain t\textgreater=0, positive constant and slope in
(33), and 0\textless n\_0\textless2U\_lL\_r/g. Both original affine
coefficient columns are replayed after EVERY original odd step and
division: the endpoint constant is odd, its slope is even, and their
divisions have exactly the declared exponent. Thus the certificate
proves the entire infinite family at each row, not only its anchor. The
other 87,943 pairs are retained incompatible labelled faces.

These whole-family height assertions are finite, computer-assisted
statements on the explicitly bounded domain (34)-(35). They are not
extended to all word pairs without proof. The general formulas (31)-(33)
remain valid for any words.

\subsection{The exact residue cover and its independent comparison}

The published one-left-edge catalogue comprises all right words of
lengths 2 through 12, with odd last exponent and
2\^{}\{A\_r-1\}\textless3\^{}\{j-1\}. The new code reconstructs its
complete 20,397 words and its stated original source formulas. The old
initial arithmetic root condition is

\SourceMath{n\bmod36\in\{3,7,15,19,27\},\quad
 n\not\equiv8731\pmod{8748},\quad n\not\equiv1731\pmod{2916}.}{36}

Before the published infinite layers are removed, the old catalogue
leaves 93,025 root residues modulo P\_0=708,588, using 370 selected
original families.

Use the common new period P=8P\_0=5,668,704. The lifted old set has
744,200 residues. Visit the compatible new pairs in the declared order
(source period, first admitted source, total word length, left word,
right word). Select a family exactly when it removes a previously
retained root residue. This yields 67 selected families and 739,770
residual residues: 4,430 are removed from this periodic set. A separate
union of ALL 24,107 compatible families, in reverse order, gives exactly
the same resulting set R\_new.

The complete 112,050-row ledger includes incompatible faces. An
independent checker imports no contribution or predecessor module,
verifies original coefficient formulas and valuations, checks strict
lexical order and the exact binomial count in EVERY domain bucket, and
independently rebuilds both residue masks. Completeness is therefore a
proved finite-domain count together with all rows, not an observed
sample. These counts describe the specified reductions; they are not
counts or probabilities of Collatz counterexamples.

The periodic count precedes the old and new infinite layer subtractions.
In particular 4,430 is not being called the number of new classes
remaining after an uncomputed overlap with the old infinite layers.

\subsection{Which of the new infinite layers lie beyond the finite
catalogue}

The new e=2 families b=1,2,3,4 and e=6 family b=1 are entirely covered
by the finite stage. Every new e=2 family b\textgreater=5 and new e=6
family b\textgreater=2 lies entirely inside its retained periodic root
set and is therefore a further removal.

For 1\textless=b\textless=10 the complete residue comparison is included
in the certificate: list every residue in the image (12) modulo P and
test the complete mask, using its exact period divisibility. At
b\textgreater=11 every source has 3\^{}11\textbar n and n=27 modulo32;
CRT gives the single residue 177147 modulo P. That residue is retained.
This proves the assertion at ALL larger layers, rather than
extrapolating from a finite b sample. Theorem 2 also proves these new
sources were not removed by any old infinite layer. Thus the cited
surviving layers supply genuinely new root reductions at arbitrarily
large valuations.

\section{6. A whole-source integral retraction, with its edge relations
retained}

Use three ordered rule stages. Stage 0 uses the published rules:
one-return forward descent; then the original image inverses of (1),
(1,2), and (1,1,1,2,1,1,4); then the old (b,e)=(1,2) join; then the
370-family catalogue; then the old infinite-layer test (22). Stage 1
uses the 67 new families after all stage 0 rules fail. Stage 2 uses
(11), with its maximal-run choice (25), after all stage 1 rules fail.
The original fixed point1 is left as a root.

Each rule gives a finite original chain B\_n and a smaller positive odd
source pi(n), with a signed integer clock k\_n and

\SourceMath{d_\epsilon B_n=V_n-q^{k_n}V_{\pi(n)}.}{37}

A root has H(V\_r)=0 and Q(V\_r)=V\_r. At any other source set

\SourceMath{H(V_n)=B_n+q^{k_n}H(V_{\pi(n)}),\quad
 Q(V_n)=q^{k_n}Q(V_{\pi(n)}),\quad F=I-Hd_\epsilon.}{38}

Every selected source decreases by at least two. The recursion therefore
terminates at a retained root for EVERY original input, without assuming
Collatz. Its intermediate original path vertices need not decrease.

\subsection{Theorem 4. Nested original chain retractions}

For each stage the maps in (38) satisfy

\SourceMath{d_\epsilon H=I-Q,\quad Q^2=Q,\quad HQ=0,
 \quad d_\epsilon F=Qd_\epsilon,\quad F^2=F,\quad FH=0.}{39}

The inclusion of {[}im F -\textgreater{} im Q{]}, its projections F,Q,
and H are an integral chain-homotopy retraction. In edge degree the
complete presentation is

\SourceMath{C^0_{\mathbb D}/\operatorname{im}H\xrightarrow{\sim}\operatorname{im}F,
 \qquad[e]\mapsto Fe.}{40}

For any earlier/later pair of the three stages, with K=H\_new-H\_old,

\SourceMath{F_{old}K=K,\quad F_{new}F_{old}=F_{old}F_{new}=F_{new},
 \quad Q_{new}Q_{old}=Q_{old}Q_{new}=Q_{new}.}{41}

\textbf{Proof.} Equation (37) telescopes to dH=I-Q. The map Q ends at
roots and H is zero there, proving Q\^{}2=Q,HQ=0. Substituting these
into F=I-Hd gives the last three identities in (39), including
FH=H-H(I-Q)=0. Also I-F=Hd is the edge-degree homotopy equation. If
Fe=0, then e=Hde is in im H; conversely FH=0 kills im H. Thus ker F=im
H, and the inverse of (40) sends Fe to {[}e{]}. In particular projected
edge columns are not assumed to be an independent new basis.

All earlier rules have priority. A later route therefore follows the
earlier route unchanged until its old root, including its terminal
q-power, and then continues there. Hence H\_new=H\_old+H\_new Q\_old and
Q\_new=Q\_new Q\_old. The later final root is also an earlier root, so
Q\_old Q\_new=Q\_new. Applying d and HQ=0 yields dK=Q\_old-Q\_new,
H\_old dK=0, hence F\_old K=K. Substitution of H\_new into both products
of I-Hd gives (41). These arguments use the actual paths and
coefficients, not only abstract isomorphism types. QED.

The choices depend on original integers, not on epsilon. Constant
reduction commutes with the maps and sends every q\^{}k to 1. It
therefore induces the same comparison-kernel equivalence (5) at every
stage. E\_1 is retained and FE\_1=E\_1. A longer actual weighted period
has boundary (1-q\^{}m)V\_c=-m epsilon V\_c; no division by m is
permitted. The integral cycle-index torsion and the free nonperiodic
summands remain represented.

For a finite declared equation support E, include its original targets
and the supports of their finite H columns. This enlargement commutes
with unions and inclusions. The resulting maps are the literal component
maps in the Split-Zero support diagram. They retain the original
relation submodule, not an inferred chronological log or an abstract
zero without its map.

\section{7. Subquadratic original-support bounds, including maximal
lifts}

Let k(n) be the largest integer k with 3\^{}k\textless=n.~Define the
integer-computable bound

\SourceMath{\boxed{R(n)=\max\{130(n+1),\lfloor64n\,4^{k(n)}/3^{k(n)}\rfloor+21\}.}}{42}

Also put

\SourceMath{\boxed{K(n)=\frac{n-1}{2}\bigl(3k(n)+14\bigr)}\qquad(n\text{ odd}).}{43}

\subsection{Theorem 5. Bounds on the complete original homotopy columns}

For every original odd n\textgreater=3, every original odd vertex in
H(V\_n), including vertices in the two uncollected paths of each
selected relation, is at most R(n). There are at most K(n) uncollected
signed monomial edge terms. The collected constant coefficients have
total absolute sum at most K(n); the epsilon coefficients have total
absolute sum at most K(n)\^{}2.

For n\textgreater=27,

\SourceMath{R(n)=\lfloor64n\,4^{k(n)}/3^{k(n)}\rfloor+21
       \le64n^{\log_3 4}+21.}{44}

Thus the original-support growth exponent is log\_3(4),
approximately1.262, rather than the earlier coarse quadratic bound. The
exact operational bound is (42), with integer powers; no floating-point
exponent is needed by the checker.

\textbf{Proof.} A selected source x in the route is at most n.~Every
finite-catalogue or inverse leg has length at most 12 and begins below
or at x. On positive inputs \nolinkurl{g_a(y)<=g_1(y)=(3y+1)/2,} and
\nolinkurl{g_1^j(y)=(3/2)^j(y+1)-1.} Thus these legs stay below
130(n+1). The three-edge left leg also obeys that bound.

For an old layer at x with b=nu\_3(x), the explicit largest initial rise
is at most 64x(4/3)\^{}b+21, and the following falling vertices do not
increase it. For a new layer, the initial rise ends at (15), at most
32x(4/3)\^{}b+21. The later vertices in (16)-(17) are bounded by this or
by (9x+7)/2, in turn below the bound with 64. The maximal-lift theorem
proves that increasing its leading run changes neither this peak nor the
later path. Since b\textless=k(n), all these values are bounded by the
second term of (42). Integer rounding gives the stated exact bound.

There are at most(n-1)/2 moves. The finite pairs have at most 15 total
edge terms. Old layers have at most 2b+14. By (30), new maximal layers
have at most \nolinkurl{A_max+b+6<=k(n)+2b+e+5<=3k(n)+11.} The uniform
factor3k(n)+14 therefore bounds all moves. This proves (43). Every term
is a signed q-power whose exponent is a sum of earlier signed leg
lengths and a position in the current leg. Its absolute value is at most
the total uncollected edge count K(n). Expansion q\^{}j=1+j epsilon and
the triangle inequality prove the two collected coefficient bounds.

For k\textgreater=3, 64(4/3)\^{}k\textgreater=4096/27\textgreater130. At
n\textgreater=27 this gap, with the retained +21, dominates130(n+1) even
after flooring. Finally3\^{}k\textless=n implies
(4/3)\^{}k\textless=(n)\^{}\{log\_3(4/3)\}, proving (44). QED.

The actual source-height maps consequently obey

\SourceMath{Q(C^1_{\le n})\subset C^1_{\le n},\quad
 H(C^1_{\le n})\subset C^0_{\le R(n)},\quad
 F(C^0_{\le n})\subset C^0_{\le R(2n)}.}{45}

The last uses T(x)\textless2x. Expanding every odd-return edge into its
original odd step and divisions gives intermediate positive integers at
most3R(n)+1 for an H column. The factors and clock assignment in that
expansion are the original ones.

These bounds apply on the WHOLE original source, including unresolved
roots. They are bounds for the chosen retraction, not bounds on all
Collatz stopping times. At a retained root its H column is zero with its
source and relations still present; no conclusion of convergence is
attached to that zero column.

\section{8. The exact remaining roots}

Let R\_new be the 739,770-residue periodic set from Section 5. The final
nonbase roots are exactly

\SourceMath{\boxed{\{n>1:n\bmod5668704\in R_{new}\}
   \setminus\left(\bigcup_{b,e}\mathcal R^{old}_{b,e}
                  \ \cup\ \bigcup_{b,e}\mathcal S_{b,e}\right).}}{46}

Both unions range over every b\textgreater=1 and e in \{2,6\}.
Membership is a finite original calculation: compute nu\_3(n), its unit,
and the matching old/new congruence. The new mask also supplies the n=27
modulo32 restriction. No infinite search over b or cutoff for maximal
depth is used. The basepoint 1 is separately retained.

Every least integer of an actual component away from 1 lies in (46): an
accepted original shared-future relation would give a smaller original
integer in that same component. This follows from the two paths, not
from assuming that a lower source is already in the base component. Root
residue counts therefore restrict possible component minima; they do not
count such components.

\section{9. The longer sharp crossing window}

The source envelope is reproduced from the preceding exact calculation,
not claimed as a new theorem in this increment. Put

\SourceMath{b_j=\operatorname{bitlength}(3^j)-1,\quad b_0=0,\quad
 C_0^*=0,\quad C_m^*=3C_{m-1}^*+2^{b_{m-1}},}{}
\SourceMath{D_m^*=2^{b_m+1}-3^m,\qquad\Theta_m=C_m^*/D_m^*.}{47}

At a first coefficient crossing with m odd returns, all proper
A\_j\textless=b\_j and the final A\_m\textgreater=b\_m+1. The exact
difference

\SourceMath{C_m^*-C_p=\sum_{j=0}^{m-1}3^{m-1-j}(2^{b_j}-2^{A_j})\ge0}{}

and D\_p\textgreater=D\_m\^{}*\textgreater0 give
C\_p/D\_p\textless=Theta\_m. Equality occurs at the original word with
proper prefix sums b\_j and final sum b\_m+1. Thus the bound is sharp on
its word domain, not necessarily attained at an integer non-descending
source. The original non-descent equation gives
n\textless=floor(Theta\_m).

For the shortened map, the first crossing with m odd steps occurs after
exactly b\_m+1 total shortened steps. Its numerator is the same C\_p,
even when it occurs before the last odd-return edge has finished all its
divisions. The envelope therefore bounds exceptions at that original,
possibly even, endpoint as well. The exact expansion of an odd exponent
a is one shortened odd step followed by a-1 shortened even steps; this
retains both clocks.

The complete integer recurrence through m=11610 gives

\SourceMath{\max_{1\le m\le11610}\lfloor\Theta_m\rfloor=8129266,
 \quad\text{record at }(m,b_m+1)=(10946,17349).}{48}

The next record is (11611,18403,9966964). All 4,064,632 original odd
sources from 3 through 8,129,265 have been scanned and independently
replayed. Each has a first coefficient crossing, strictly below its
source at BOTH the first shortened crossing and the final odd-return
endpoint. There are 14,186,680 odd-return relations and 28,378,002
independently repeated divisions in the full audit.

Together with the proved envelope, this finite certificate establishes:

\begin{quote}
Every positive integer n\textgreater=2 whose first shortened coefficient
crossing involves at most 11,610 odd steps strictly descends at that
crossing.
\end{quote}

Even starting integers descend at the first division; the odd cases
above and the envelope prove the rest. The starting integer and final
exponent are not bounded in this statement. Sources with no such
crossing remain unclaimed.

The two producer runs, ordinary and optimized, are byte-identical. The
independent auditor imports no contribution or predecessor module, reads
EVERY source row, performs the original divisions separately,
reconstructs the envelope with a critical-prefix recurrence, and checks
the smaller-source induction table. The largest observed first crossing
is 155 odd returns at 8,088,063; the largest inductive total is 248 odd
returns at 6,649,279. Those observations are not used as unproved global
clock bounds. The largest odd vertex appearing in the audited crossing
paths is 20,114,203,639,877, reached from 6,631,675; all large values
are retained as exact integers.

The smaller positive odd endpoints supply strong induction to 1 on the
entire finite interval. In particular the height-filtered original
kernel F\_N=\textless xi\_n:n\textless=N\textgreater{} has F\_8129265=0,
with its original finite integral witnesses. An arbitrary lower-source
relation xi\_n=xi\_m with m\textless n only kills the associated height
increment; it does not erase that retained lower subgroup.

Coefficient stopping time and merging/sufficient-set methods have
established literature; see {[}RT{]} and {[}M{]}. No priority or
published verification record is claimed. The new unrestricted family
proofs do not depend on the scan or on a literature bound.

\section{10. What the results do and do not close}

Any least vertex of a nonbase component must exceed 8,129,265, lie in
(46), and have no first coefficient crossing involving at most 11,610
odd steps. These are simultaneous necessary conditions on original
positive integers. The proof has not shown that their intersection is
empty.

The new all-layer domains are disjoint from the published ones, and
their maximal leading runs now have complete residue strata and bounds.
The finite two-sided domain is fully certified, and its exact cover is
explicit. The whole-system retraction has a smaller global
support-growth bound while keeping its residual component module. None
of these facts assumes that a finite open frontier is an infinite orbit
or that a zero in a rationalized module vanishes integrally.

The next arithmetic targets are further right-leg templates or genuinely
new left-leg patterns on (46), and residue-aware restrictions on the
exact canonical crossing sources. The next fixed crossing-record
calculation has horizon 11611 and odd candidate ceiling 9,966,963. These
are stated remaining calculations, not scheduled tasks or assumed
results. A single ever-larger finite certificate does not establish that
every original singleton vanishes.

The predecessor bound Theta\_m\textgreater m/6 follows from
2\^{}\{b\_j\}\textgreater3\^{}j/2 and D\_m\^{}*\textless3\^{}m, so no
fixed finite source interval covers all crossing horizons under that
envelope. A global proof needs additional arithmetic on the original
remaining source, or a rigorously surviving integral source class for a
counterexample. The supplied maps preserve that distinction.

\section{11. Executed evidence and source lineage}

\nolinkurl{verify.py} checks the complete 112,050-pair domain, all
24,107 compatible variable families, all 20,397 old candidate families,
both residue unions, the layer formulas and saturation strata, original
signed chain maps, the nested retractions, support/length bounds, and
four original singleton witness extractions. It rejects seventeen
malformed controls. General family statements have the written proofs
above; finite symbolic-family rows are complete finite proof components,
and point fixtures are regression evidence only.

The structural pair has 1,905,007 named checks per run and
byte-identical output. The all-pair auditor independently checks 112,050
rows and 332,467 original affine row equations. The full source scan and
its independent audit have the separate scopes in Section 9. Counts are
not summed and reported as a number of theorems. The actual command
receipts and file hashes identify the executed bytes.

The source programme and Split-Zero framework are attributed to
KokunoYumeto. This tool-assisted continuation develops the new
arithmetic, proofs, and checks; it is not an independent external
mathematical review. No new Lean execution is claimed. The complete
earlier contribution is neither relicensed nor rewritten.

{[}WB{]} KokunoYumeto, Collatz workbench,
\nolinkurl{ternary_join_reduction_20260916/note.md}, commit
\nolinkurl{b2dac4c2cc4a5eedde0af36d6749390c77fa4a62.} Full relevant
source text and implementation priority were read through the
authenticated GitHub connector. The new code reconstructs its arithmetic
rules; it does not claim that a full remote checkout was replayed. The
local preceding height-window source is also retained as lineage, not
silently reported as newly published.

{[}IJ{]} \nolinkurl{intrinsic_zero_firstjet_20260915/firstjet.py}, Git
blob \nolinkurl{97f9ef5ad535f7ec18259d41fab6a3557be17d98,} is the
unchanged executable receiver. Its byte identity is checked before every
imported witness verification.

{[}RT{]} Olivier Rozier and Claude Terracol, \emph{Paradoxical behavior
in Collatz sequences}, arXiv:2502.00948v5 (17 May 2026), definitions
1.1--1.2 and discussion of Terras' coefficient stopping-time conjecture.
Source for attribution, not an input to the present arithmetic proofs.
\nolinkurl{https://arxiv.org/html/2502.00948v5}

{[}M{]} Keenan Monks, Kenneth G. Monks, Kenneth M. Monks and Maria
Monks, \emph{Strongly sufficient sets and the distribution of arithmetic
sequences in the 3x+1 graph}, arXiv:1204.3904. Prior
merging/sufficient-set context; no result from that paper is used
without an additional map or invoked as a premise here.
\nolinkurl{https://arxiv.org/abs/1204.3904}
